% ======================================================================
% capitulos/capitulo1.tex
% ----------------------------------------------------------------------
% ENTREGABLE 1 — Diagnóstico Organizacional y Alineamiento Estratégico
% Empresa caso de estudio: WISP Ingeniería en Redes y Telecomunicaciones
% Contenido trasladado desde Entregable1_GTI_WISP.docx
% ======================================================================
\chapter{Diagnóstico Organizacional y Alineamiento Estratégico}
\label{chap:diagnostico}

% ----------------------------------------------------------------------
% Descripción del producto (introducción del capítulo)
% ----------------------------------------------------------------------
El presente documento constituye el Entregable~1 de la línea de Gestión de
Tecnologías de Información (GTI) del Proyecto de Evaluación de Fin de Carrera.
Corresponde a un análisis estructurado de la situación actual de la empresa
WISP Ingeniería en Redes y Telecomunicaciones, orientado a identificar sus
brechas estratégicas, operativas y tecnológicas.

El análisis incluye la evaluación del contexto interno y externo de la
organización, el grado de madurez en tecnologías de información y la coherencia
entre los objetivos estratégicos institucionales y las capacidades digitales
existentes. Este componente fundamenta la necesidad del proyecto y garantiza su
alineamiento con el plan estratégico organizacional y el plan estratégico de TI,
sentando las bases para el desarrollo de una solución tecnológica integral que
abarca las líneas de Gestión de TI, Ingeniería de Software e Infraestructura
Tecnológica.

\begin{upeubox}{Alcance del documento}
Este entregable cubre los cuatro componentes exigidos por la rúbrica:
(1.1)~Contexto Organizacional, (1.2)~Análisis Estratégico,
(1.3)~Diagnóstico Digital~/~TI y (1.4)~Identificación del Problema. Los
resultados de este diagnóstico alimentan directamente el Business Case, el Plan
de Gestión del Proyecto y el diseño de la solución técnica integrada.
\end{upeubox}

% ======================================================================
\section{Contexto Organizacional}
\label{sec:contexto}

\subsection{Descripción del sector y entorno competitivo}

WISP Ingeniería en Redes y Telecomunicaciones (en adelante, WISP) es una empresa
peruana del sector de telecomunicaciones, especializada en la provisión del
servicio de acceso a internet mediante red de fibra óptica (FTTH~--~\textit{Fiber
To The Home}) y televisión digital en zonas rurales y semi-rurales de la región
Puno. La empresa opera en un segmento de mercado caracterizado por una baja
penetración de servicios de conectividad de banda ancha y una limitada presencia
de los operadores de telecomunicaciones de escala nacional.

El sector de telecomunicaciones en el Perú es un mercado regulado por el
Organismo Supervisor de Inversión Privada en Telecomunicaciones (OSIPTEL) y el
Ministerio de Transportes y Comunicaciones (MTC). En las grandes ciudades de la
región, como Puno y Juliaca, la competencia es intensa y está conformada tanto
por operadores nacionales ---Movistar, Claro, Entel y WOW--- como por operadores
locales de fibra óptica de creciente presencia, entre los que destacan Fibramax,
WiwiTel, FiberUltra y otros proveedores regionales. Estos actores concentran su
inversión en las zonas urbanas de alta densidad poblacional, donde la
rentabilidad por kilómetro de red desplegada es mayor.

Sin embargo, en los distritos rurales alejados donde opera WISP, la presencia de
estos operadores es marginal o inexistente, debido a la baja densidad
poblacional, la dispersión geográfica y la complejidad logística que implica el
despliegue de infraestructura de fibra en dichas zonas. Esta característica del
mercado configura una ventaja competitiva estructural para WISP: la empresa
atiende un mercado desatendido por los grandes operadores y por la mayoría de los
ISP locales, contribuyendo activamente a la reducción de la brecha digital en la
región Puno. WISP ya cuenta con infraestructura de fibra óptica desplegada y
operativa en trece~(13) zonas distribuidas en distintas provincias de la región.

\subsubsection{Zonas de cobertura y modalidad de servicio}

WISP brinda el servicio de internet por fibra óptica en las siguientes zonas de
la región Puno:

\begin{center}
\captionof{table}{Zonas de cobertura de fibra óptica de WISP en la región Puno.}
\label{tab:zonas}
\begin{tabularx}{\textwidth}{C{1.1cm} X C{1.1cm} X}
\toprule
\textbf{N°} & \textbf{Zona / Distrito} & \textbf{N°} & \textbf{Zona / Distrito} \\
\midrule
1 & Cabana        & 8  & Mañazo  \\
2 & Cabanilla     & 9  & Muñani  \\
3 & Cabanillas    & 10 & Ñuñoa   \\
4 & Capachica     & 11 & Pucara  \\
5 & Choquehuanca  & 12 & Sucasco \\
6 & Domingo       & 13 & Vilque  \\
7 & Huata         &    &         \\
\bottomrule
\end{tabularx}
\end{center}

\begin{upeubox}{Nota sobre el entorno competitivo (dato de campo)}
Debe distinguirse entre el mercado urbano y el rural, porque la competencia que
enfrenta WISP es radicalmente distinta en cada uno:
\begin{itemize}[leftmargin=1.4em]
  \item \textbf{En las ciudades (Puno y Juliaca):} el mercado es altamente
  competitivo, con operadores nacionales (Movistar, Claro, Entel, WOW) y locales
  de fibra (Fibramax, WiwiTel, FiberUltra, entre otros). WISP no compite
  directamente en este segmento.
  \item \textbf{En los distritos rurales (zona de operación de WISP):} el único
  competidor con fibra óptica es FiberUltra, y solo en dos distritos: Cabanillas
  y Mañazo. En el resto de zonas ---e incluso en esos dos distritos--- la
  competencia se limita a pequeños proveedores de internet inalámbrico
  (radioenlaces), de menor estabilidad y capacidad que la fibra. En la mayoría de
  sus zonas, WISP es el único proveedor de fibra óptica, una posición de ventaja
  competitiva difícil de replicar en el corto plazo.
\end{itemize}
\end{upeubox}

\subsection{Estructura organizacional}

WISP opera bajo una estructura organizacional funcional y de tamaño reducido,
característica de una micro/pequeña empresa (MYPE) del sector tecnológico. La toma
de decisiones está centralizada en la administración, y las áreas operativas
(técnica y comercial) se despliegan sobre las trece zonas de cobertura. La
estructura actual del personal es la siguiente:

\begin{center}
\includegraphics[width=\textwidth]{organigrama.png}
\captionof{figure}{Estructura del personal de WISP por área y responsabilidades.}
\label{fig:estructura}
\end{center}

\begin{upeubox}{Observación estructural crítica}
Un solo administrador concentra la gestión de aproximadamente 455 clientes
distribuidos en 13 zonas geográficamente dispersas. Esta alta concentración de
responsabilidades en una sola persona constituye un \textbf{punto único de falla}
operativo y una limitación severa para el crecimiento, aspecto que se profundiza
en el diagnóstico de TI (Sección~\ref{sec:diagnostico}) y en la identificación
del problema (Sección~\ref{sec:problema}).
\end{upeubox}

\subsection{Cadena de valor}

Siguiendo el modelo de Michael Porter, la cadena de valor de WISP descompone la
empresa en sus actividades constitutivas para identificar las fuentes de ventaja
competitiva. Las actividades se dividen en primarias (directamente involucradas
en la creación y entrega del servicio al cliente) y de apoyo (que respaldan a las
primarias). El margen resulta de la diferencia entre el valor que el cliente está
dispuesto a pagar y el costo de desarrollar dichas actividades.

El análisis revela un hallazgo central: el flujo de información atraviesa
transversalmente toda la cadena, y su gestión deficiente ---actualmente basada en
WhatsApp, papel y un sistema heredado--- erosiona el margen al elevar los costos
operativos y de coordinación en cada eslabón.

\begin{center}
\includegraphics[width=\textwidth]{cadena_de_valor.png}
\captionof{figure}{Cadena de valor de WISP según el modelo de Michael Porter.}
\label{fig:cadena-valor}
\end{center}

\begin{upeubox}{Diagnóstico de la cadena de valor}
La ventaja competitiva de WISP se sustenta hoy en su Desarrollo Tecnológico de red
(infraestructura de fibra madura) y en su cercanía al cliente en Marketing/Ventas
y Servicio postventa. Sin embargo, el margen se ve comprometido porque la
Infraestructura de la empresa (sistemas de información y gestión) opera en un nivel
deficiente que encarece y ralentiza todas las actividades primarias. El proyecto
interviene precisamente sobre esta actividad de apoyo transversal para liberar el
potencial de toda la cadena.
\end{upeubox}

\subsection{Mapa de stakeholders}

Se identifican los grupos de interés (\textit{stakeholders}) internos y externos
que se relacionan con la operación de WISP y que serán impactados directa o
indirectamente por el proyecto de transformación tecnológica. Se clasifican según
su naturaleza, su rol y su nivel de influencia e interés en el proyecto.

\begin{center}
\captionof{table}{Mapa de stakeholders del proyecto.}
\label{tab:stakeholders}
\begin{tabularx}{\textwidth}{L{2.9cm} C{1.4cm} X C{1.5cm} C{1.3cm}}
\toprule
\textbf{Stakeholder} & \textbf{Tipo} & \textbf{Rol / Interés principal} & \textbf{Influencia} & \textbf{Interés} \\
\midrule
Administrador / Gerente & Interno &
Decisor. Busca control total de la operación y eficiencia administrativa. & Alta & Alto \\
Técnicos (4) & Interno &
Operativo. Requieren herramientas que simplifiquen instalaciones y registro de averías. & Media & Alto \\
Vendedoras tienda (2) & Interno &
Comercial. Registro rápido y validación de prospectos. & Media & Alto \\
Vendedoras campo (4) & Interno &
Comercial. Necesitan app móvil para registrar clientes sin papel ni WhatsApp. & Media & Alto \\
Clientes finales ($\approx$455) & Externo &
Usuario. Buscan servicio estable, atención rápida y comunicación clara. & Baja & Alto \\
Proveedor mayorista de internet & Externo &
Soporte. Provee el transporte de última milla; interesa continuidad. & Alta & Medio \\
Competidores (FiberUltra, ISP inalámbricos) & Externo &
Amenaza. Compiten por captación de clientes en zonas compartidas. & Media & Bajo \\
Entes reguladores (OSIPTEL, MTC) & Externo &
Regulación. Supervisan cumplimiento normativo del servicio. & Alta & Bajo \\
\bottomrule
\end{tabularx}
\end{center}

\begin{upeubox}{Priorización de stakeholders (matriz Poder--Interés)}
Los stakeholders de mayor prioridad para el proyecto (alta influencia $+$ alto
interés) son el administrador y, en su conjunto, el personal operativo (técnicos y
vendedoras). Estos serán los usuarios directos del sistema y su adopción
determinará el éxito de la iniciativa. Los clientes finales, aunque de baja
influencia directa, tienen alto interés porque percibirán la mejora en la calidad
y rapidez del servicio.
\end{upeubox}

% ======================================================================
\section{Análisis Estratégico}
\label{sec:estrategico}

El análisis estratégico establece el marco de referencia sobre el cual se sustenta
la necesidad del proyecto. Se formaliza la misión, visión y objetivos estratégicos
de WISP, se evalúa su situación interna y externa mediante el análisis FODA, se
complementa con un análisis PESTEL del macroentorno y se determinan los factores
críticos de éxito.

\subsection{Misión, visión y objetivos estratégicos}

\begin{upeubox}{Justificación de la propuesta}
Dado que WISP, por ser una MYPE de estructura reducida y decisiones
centralizadas, no contaba con una misión y visión formalmente declaradas,
este documento las formula por primera vez aplicando el marco teórico del
planeamiento estratégico.
\end{upeubox}

\begin{upeubox}{Misión (propuesta formalizada)}
"Proveer acceso a internet de alta velocidad por fibra óptica y televisión digital,
con estándares de calidad, estabilidad y trazabilidad, a las familias y negocios
de las zonas de la región Puno, sustentados en infraestructura propia,
un servicio técnico cercano y un modelo de gestión eficiente que garantiza
continuidad y precios accesibles."
\end{upeubox}

La \textbf{misión} responde a las tres preguntas fundamentales del modelo de
formulación estratégica: \emph{qué} necesidad satisface WISP (acceso a internet
de calidad en zonas desatendidas), \emph{cómo} la satisface (mediante
infraestructura propia de fibra óptica y atención cercana) y \emph{a quiénes}
atiende (las familias y negocios de las comunidades de la región Puno). Su
contenido no es arbitrario: se deriva del diagnóstico organizacional (FODA,
cadena de valor) y de los valores institucionales de la empresa.

\begin{upeubox}{Visión (propuesta formalizada)}
"Al 2030, ser el proveedor líder de internet por fibra óptica de Puno,
llevando conectividad de calidad y nuevas oportunidades a las zonas donde otros no llegan."
\end{upeubox}

La \textbf{visión} cumple las cinco características que exige el marco teórico:
es \emph{inspiradora} (proyecta el propósito de llegar donde otros no llegan),
\emph{clara y concisa} (con un horizonte definido al 2030), \emph{ambiciosa
pero realista} (el liderazgo en fibra es una posición que WISP ya casi
ocupa), \emph{relevante y sostenible} (alineada con los objetivos estratégicos
de expansión y escalabilidad) y \emph{flexible y adaptativa} (no se limita a un
número fijo de zonas, permitiendo el crecimiento futuro).

En conjunto, misión y visión quedan alineadas entre sí y con los valores
institucionales, garantizando la coherencia estratégica que el planeamiento
requiere.

\subsubsection{Valores institucionales}

Los valores constituyen la brújula ética que orienta el comportamiento y las
decisiones de WISP, y sirven de base a su cultura organizacional:

\begin{center}
\captionof{table}{Valores institucionales de WISP.}
\label{tab:valores}
\begin{tabularx}{\textwidth}{L{3.4cm} X}
\toprule
\textbf{Valor} & \textbf{Significado en WISP} \\
\midrule
Cercanía & Atención personalizada y trato directo con el cliente rural, entendiendo su realidad y necesidades. \\
Compromiso social & Vocación de cerrar la brecha digital y llevar conectividad a zonas que otros operadores no atienden. \\
Calidad y confiabilidad & Provisión de un servicio estable, con infraestructura de fibra óptica robusta y soporte técnico propio. \\
Innovación & Adopción de tecnología para mejorar continuamente la operación y la experiencia del cliente. \\
Integridad & Transparencia en el cobro, en los contratos y en la relación con clientes y colaboradores. \\
\bottomrule
\end{tabularx}
\end{center}

A partir de la misión, visión y valores, se derivan los siguientes objetivos
estratégicos institucionales:

\begin{center}
\captionof{table}{Objetivos estratégicos institucionales de WISP.}
\label{tab:objetivos}
\begin{tabularx}{\textwidth}{C{1.4cm} X L{2.3cm}}
\toprule
\textbf{Código} & \textbf{Objetivo estratégico} & \textbf{Perspectiva} \\
\midrule
OE-01 & Expandir la cobertura de fibra óptica a nuevos distritos de la región Puno actualmente desatendidos. & Crecimiento \\
OE-02 & Mejorar la eficiencia operativa mediante la digitalización y automatización de los procesos de captación, instalación, soporte y gestión de clientes. & Procesos \\
OE-03 & Incrementar la satisfacción y retención de clientes a través de un servicio estable y una atención más rápida y trazable. & Cliente \\
OE-04 & Fortalecer el control administrativo y financiero, reduciendo la morosidad y la dependencia de procesos manuales. & Financiera \\
OE-05 & Consolidar una plataforma tecnológica escalable que soporte el crecimiento de la empresa sin incrementar proporcionalmente la carga administrativa. & Tecnología \\
\bottomrule
\end{tabularx}
\end{center}

\subsection{Análisis FODA}

El análisis FODA sintetiza los factores internos (fortalezas y debilidades) y
externos (oportunidades y amenazas) que condicionan la posición estratégica de
WISP. Se presenta primero en forma de matriz y luego se detalla cada cuadrante.

\begin{center}
\captionof{table}{Matriz FODA de WISP.}
\label{tab:foda}
\begin{tabularx}{\textwidth}{X X}
\toprule
\rowcolor{upeublue!10}
\textbf{FORTALEZAS (F)} & \textbf{DEBILIDADES (D)} \\
\midrule
\begin{itemize}[leftmargin=1.4em,nosep]
  \item F1. Infraestructura propia de fibra óptica en 13 zonas.
  \item F2. Único proveedor de fibra en la mayoría de sus zonas.
  \item F3. Servicio técnico propio (4 técnicos).
  \item F4. Fuerza de ventas en campo desplegada por zonas.
  \item F5. Portafolio de planes con TV integrada (\textit{triple play}).
  \item F6. Conocimiento del mercado rural y cercanía con el cliente.
\end{itemize}
&
\begin{itemize}[leftmargin=1.4em,nosep]
  \item D1. Sistema de gestión heredado (PHP) lento y mal diseñado.
  \item D2. Operación dependiente de WhatsApp, sin trazabilidad.
  \item D3. Sin aplicación móvil para técnicos ni vendedoras.
  \item D4. Un solo administrador para 13 zonas (punto único de falla).
  \item D5. Sin historial técnico por cliente (caja NAP, borne, ONU).
  \item D6. Inconsistencia de datos entre MikroTiks y planes comerciales.
\end{itemize}
\\
\midrule
\rowcolor{upeublue!10}
\textbf{OPORTUNIDADES (O)} & \textbf{AMENAZAS (A)} \\
\midrule
\begin{itemize}[leftmargin=1.4em,nosep]
  \item O1. Numerosos distritos rurales aún sin cobertura de fibra.
  \item O2. Baja competencia de fibra (solo inalámbricos en la mayoría de zonas).
  \item O3. Programas estatales de conectividad rural y masificación de internet.
  \item O4. Creciente demanda de internet (educación, trabajo, streaming).
  \item O5. Abaratamiento de tecnología GPON y equipos de red.
\end{itemize}
&
\begin{itemize}[leftmargin=1.4em,nosep]
  \item A1. Competidor FiberUltra con fibra en Cabanillas y Mañazo.
  \item A2. Proveedores inalámbricos que compiten por precio.
  \item A3. Posible expansión futura de operadores nacionales a zonas rurales.
  \item A4. Cambios regulatorios de OSIPTEL/MTC.
  \item A5. Dependencia de un proveedor mayorista de internet.
\end{itemize}
\\
\bottomrule
\end{tabularx}
\end{center}

\subsubsection{Estrategias derivadas del cruce FODA}

Del cruce entre los factores internos y externos se derivan estrategias que
orientan la intervención del proyecto:

\begin{center}
\captionof{table}{Estrategias derivadas del cruce FODA.}
\label{tab:foda-cruce}
\begin{tabularx}{\textwidth}{L{2.7cm} L{2.5cm} X}
\toprule
\textbf{Estrategia} & \textbf{Cruce} & \textbf{Descripción} \\
\midrule
Ofensiva (FO) & F1$\cdot$F2 $+$ O1$\cdot$O2 &
Aprovechar la infraestructura propia y la baja competencia para expandirse
agresivamente a nuevos distritos rurales sin fibra. \\
Adaptativa (DO) & D1$\cdot$D2$\cdot$D3 $+$ O4 &
Digitalizar la operación (sistema web $+$ apps móviles) para atender la creciente
demanda con eficiencia y escalar sin colapsar la administración.
\textrightarrow{}~\textit{Justifica directamente el proyecto.} \\
Defensiva (FA) & F3$\cdot$F6 $+$ A1$\cdot$A2 &
Reforzar la ventaja de servicio técnico propio y cercanía al cliente para
diferenciarse de FiberUltra y de los inalámbricos. \\
Supervivencia (DA) & D4$\cdot$D5 $+$ A1$\cdot$A3 &
Superar la dependencia de una sola persona y la falta de trazabilidad antes de que
la competencia o la escala pongan en riesgo la continuidad del negocio. \\
\bottomrule
\end{tabularx}
\end{center}

\subsection{Análisis PESTEL del macroentorno}

El análisis PESTEL complementa el diagnóstico externo evaluando las fuerzas del
macroentorno que afectan a WISP:

\begin{center}
\captionof{table}{Análisis PESTEL del macroentorno de WISP.}
\label{tab:pestel}
\begin{tabularx}{\textwidth}{L{2cm} X L{2.3cm}}
\toprule
\textbf{Factor} & \textbf{Descripción e impacto en WISP} & \textbf{Efecto} \\
\midrule
Político &
Políticas públicas de cierre de brecha digital y programas de conectividad rural
(FITEL/PRONATEL) favorecen la expansión del servicio. & Oportunidad \\
Económico &
Economía rural con ingresos limitados exige planes accesibles; a la vez, el
internet se ha vuelto un servicio básico y su demanda es inelástica. & Mixto \\
Social &
Aumento del teletrabajo, educación virtual y consumo de streaming eleva la demanda
y la exigencia de estabilidad del servicio. & Oportunidad \\
Tecnológico &
Madurez y abaratamiento de GPON, MikroTik y cloud permiten operar y escalar a bajo
costo; RouterOS~7 habilita integración por API. & Oportunidad \\
Ecológico &
Presión por eficiencia energética y correcta disposición de residuos electrónicos
(equipos, ONU en desuso). & Neutro \\
Legal &
Regulación de OSIPTEL/MTC sobre calidad de servicio, contratos y protección de
datos personales (Ley N° 29733). & Amenaza \\
\bottomrule
\end{tabularx}
\end{center}

\subsection{Análisis de las cinco fuerzas de Porter}

El modelo de las cinco fuerzas de Porter permite evaluar el nivel de competencia e
intensidad del sector en el que opera WISP, y determinar el atractivo del mercado.
Se analiza cada fuerza en el contexto específico de las zonas rurales de Puno donde
la empresa presta servicio.

\begin{center}
\captionof{table}{Análisis de las cinco fuerzas de Porter aplicado a WISP.}
\label{tab:porter}
\begin{tabularx}{\textwidth}{L{3.2cm} X L{2cm}}
\toprule
\textbf{Fuerza competitiva} & \textbf{Análisis aplicado a WISP} & \textbf{Intensidad} \\
\midrule
1. Amenaza de nuevos competidores &
El despliegue de fibra en zonas rurales exige alta inversión, conocimiento técnico y
logística compleja, lo que genera fuertes barreras de entrada. Sin embargo, los ISP
inalámbricos entran con bajo costo. & Media \\
2. Poder de negociación de proveedores &
WISP depende de un proveedor mayorista de transporte de internet (última milla). La
concentración de esta oferta otorga poder al proveedor. Los equipos (MikroTik, OLT,
ONU) tienen múltiples proveedores. & Alta \\
3. Poder de negociación de clientes &
En las zonas donde WISP es el único proveedor de fibra, el poder del cliente es
bajo. Donde compite con inalámbricos o FiberUltra, el cliente puede comparar precio
y calidad, elevando su poder. & Baja--Media \\
4. Amenaza de productos sustitutos &
El internet móvil (4G/5G) de operadores nacionales y el internet inalámbrico local
son sustitutos parciales, de menor estabilidad, pero disponibles. El satelital (ej.
Starlink) es un sustituto emergente de mayor costo. & Media \\
5. Rivalidad entre competidores &
En el segmento rural la rivalidad es baja: FiberUltra solo en 2 distritos y algunos
inalámbricos. En Puno/Juliaca la rivalidad es alta (Movistar, Claro, Entel, WOW,
Fibramax, WiwiTel), pero WISP no compite allí. & Baja (rural) \\
\bottomrule
\end{tabularx}
\end{center}

\begin{upeubox}{Conclusión del análisis de Porter}
El sector rural donde opera WISP resulta atractivo: la rivalidad directa y el poder
de los clientes son bajos gracias a la posición de único proveedor de fibra en la
mayoría de zonas, y las barreras de entrada son significativas. La principal presión
proviene del poder del proveedor mayorista de internet y de la amenaza de sustitutos
(internet móvil, satelital e inalámbrico). La estrategia de WISP debe orientarse a
consolidar su ventaja de fibra mediante eficiencia operativa y expansión, que es
justamente lo que habilita el proyecto.
\end{upeubox}

\subsection{Factores críticos de éxito (FCE)}

Los factores críticos de éxito son aquellas condiciones que WISP debe garantizar
obligatoriamente para cumplir sus objetivos estratégicos y para que el proyecto de
transformación tecnológica genere el valor esperado:

\begin{description}[leftmargin=0pt, style=nextline]
  \item[FCE-1. Estabilidad y calidad del servicio:] mantener una red confiable con
  mínima interrupción, base de la retención de clientes en un mercado donde el
  internet es servicio esencial.
  \item[FCE-2. Digitalización de la operación:] contar con un sistema centralizado
  que reemplace WhatsApp y el sistema heredado, con trazabilidad de extremo a extremo.
  \item[FCE-3. Escalabilidad administrativa:] poder crecer en clientes y zonas sin
  aumentar proporcionalmente la carga sobre la administración.
  \item[FCE-4. Agilidad en captación e instalación:] reducir el tiempo entre que un
  vendedor capta un cliente y el técnico lo activa.
  \item[FCE-5. Control financiero y de morosidad:] gestionar cortes y reactivaciones
  de forma oportuna y automatizada por zona.
  \item[FCE-6. Adopción por el personal:] que técnicos y vendedoras usen efectivamente
  las herramientas digitales desde su celular.
\end{description}

% ======================================================================
\section{Diagnóstico Digital / TI}
\label{sec:diagnostico}

Esta sección evalúa el estado actual de las capacidades digitales y tecnológicas de
WISP. Se inventarían los sistemas de información existentes, se determina el nivel de
madurez digital de la organización, se identifican las brechas tecnológicas y se
detallan los problemas operativos asociados a la gestión de TI.

\subsection{Inventario de sistemas de información existentes}

El siguiente inventario recoge las herramientas y sistemas que WISP utiliza
actualmente para soportar su operación:

\begin{center}
\captionof{table}{Inventario de sistemas de información existentes en WISP.}
\label{tab:inventario}
\begin{tabularx}{\textwidth}{L{2.9cm} X L{2.7cm} L{1.7cm}}
\toprule
\textbf{Sistema / Herramienta} & \textbf{Función} & \textbf{Tecnología} & \textbf{Estado} \\
\midrule
Sistema heredado de clientes &
Registro y administración de clientes, planes y cobros. &
PHP $+$ frontend Bootstrap (web) & Deficiente \\
MikroTik RouterOS 7 ($\times$13) &
Gestión de usuarios PPPoE, control de ancho de banda y corte de morosos por zona. &
RouterOS 7.13.5 / RB4011iGS+ & Operativo \\
OLT VSOL ($\times$13) &
Concentración de fibra óptica y gestión de ONU/ONT por zona. &
GPON / VSOL & Operativo \\
Bot de pagos WhatsApp &
Recepción de comprobantes de pago (Yape, Plin, transferencia). &
Node.js (WhatsApp Web) & Parcial \\
WhatsApp operativo &
Canal informal para captación, coordinación de instalaciones y envío de evidencias. &
WhatsApp (mensajería) & Informal \\
Registros manuales / papel &
Anotación de prospectos en campo por las vendedoras. &
Papel / hojas sueltas & Deficiente \\
\bottomrule
\end{tabularx}
\end{center}

\subsection{Nivel de madurez digital}

Para evaluar la madurez digital de WISP se utiliza una escala de cinco niveles
(Inicial, Básico, Definido, Gestionado, Optimizado), aplicada a distintas
dimensiones de la organización. El resultado ubica a WISP predominantemente en un
nivel Inicial--Básico.

\begin{center}
\captionof{table}{Nivel de madurez digital de WISP por dimensión.}
\label{tab:madurez}
\begin{tabularx}{\textwidth}{L{3cm} L{2.2cm} X}
\toprule
\textbf{Dimensión} & \textbf{Nivel actual} & \textbf{Situación observada} \\
\midrule
Sistemas de información & Inicial &
Sistema heredado deficiente y herramientas dispersas sin integración. \\
Automatización de procesos & Inicial &
Procesos mayormente manuales; WhatsApp como eje operativo. \\
Gestión de datos & Inicial &
Datos dispersos entre MikroTiks, papel y chats; sin fuente única de verdad. \\
Infraestructura de red & Gestionado &
Red de fibra y equipos MikroTik/OLT sólidos y estandarizados. \\
Movilidad / trabajo en campo & Inicial &
Sin apps móviles; el personal usa navegador y WhatsApp desde el celular. \\
Experiencia del cliente & Básico &
Canal de pagos por WhatsApp, pero sin autoservicio ni información de cobertura en línea. \\
\bottomrule
\end{tabularx}
\end{center}

\begin{upeubox}{Lectura del nivel de madurez}
WISP presenta una notable asimetría: su infraestructura de red (capa física) está
madura y bien gestionada, pero su capa de sistemas de información, datos y
automatización se encuentra en un nivel inicial. Esta brecha entre una red robusta y
una gestión digital incipiente es precisamente el espacio de oportunidad del proyecto.
\end{upeubox}

\subsection{Brechas tecnológicas}

Se identifican las siguientes brechas entre la situación actual (AS-IS) y la
situación deseada (TO-BE) que el proyecto busca cerrar:

\begin{center}
\captionof{table}{Brechas tecnológicas entre la situación actual y la deseada.}
\label{tab:brechas}
\begin{tabularx}{\textwidth}{L{2.9cm} X X}
\toprule
\textbf{Brecha} & \textbf{Situación actual} & \textbf{Situación deseada} \\
\midrule
B1. Integración de sistemas &
Sistema PHP, MikroTiks, papel y WhatsApp funcionan de forma aislada. &
Plataforma única (web $+$ móvil) integrada con los 13 MikroTiks vía API. \\
B2. Movilidad &
Sin apps; técnicos y vendedoras dependen de WhatsApp y navegador lento. &
App móvil con roles de técnico y vendedor para registro en campo. \\
B3. Trazabilidad &
No hay historial técnico ni seguimiento del ciclo de vida del cliente. &
Trazabilidad completa: prospecto \textrightarrow{} validación \textrightarrow{} instalación \textrightarrow{} activación. \\
B4. Automatización de red &
Alta/corte/reactivación de PPPoE se hace manualmente en cada MikroTik. &
Gestión automática de PPPoE y cortes desde el sistema central. \\
B5. Consistencia de datos &
Datos de cliente embebidos en comentarios del MikroTik; planes inconsistentes. &
Base de datos centralizada como fuente única de verdad. \\
B6. Escalabilidad &
La carga administrativa crece linealmente con cada nuevo cliente/zona. &
Arquitectura en nube escalable que soporta el crecimiento. \\
\bottomrule
\end{tabularx}
\end{center}

\subsection{Problemas operativos asociados a TI}

Como consecuencia de las brechas anteriores, WISP enfrenta los siguientes problemas
operativos concretos en su día a día:

\begin{description}[leftmargin=0pt, style=nextline]
  \item[Captación ineficiente:] las vendedoras de campo anotan datos en papel o los
  envían por WhatsApp, lo que genera pérdida de información, errores de transcripción
  y demoras en la validación.
  \item[Instalaciones sin registro estructurado:] los técnicos envían por WhatsApp las
  fotos de la caja NAP, el borne del splitter, los equipos instalados y la fachada, sin
  vincularlas formalmente al cliente ni conservar un historial técnico.
  \item[Gestión manual de la red:] el administrador debe ingresar manualmente a cada
  MikroTik para dar de alta, cortar o reactivar clientes PPPoE, tarea repetitiva y
  propensa a errores en 13 zonas.
  \item[Información dispersa y no confiable:] no existe una fuente única de datos; la
  información del cliente vive parcialmente en el sistema PHP, en comentarios del
  MikroTik y en chats de WhatsApp.
  \item[Dependencia de una sola persona:] toda la operación administrativa recae en el
  administrador, generando cuellos de botella y riesgo de continuidad.
  \item[Sistema heredado lento:] el frontend actual en Bootstrap es lento y poco usable
  desde el celular, dificultando el trabajo de quienes no tienen acceso a una computadora.
\end{description}

% ======================================================================
\section{Identificación del Problema}
\label{sec:problema}

\subsection{Definición estructurada del problema}

\begin{upeubox}{Problema central}
WISP Ingeniería en Redes y Telecomunicaciones gestiona la operación de
aproximadamente 455 clientes de internet por fibra óptica, distribuidos en 13 zonas
geográficamente dispersas, mediante un sistema heredado deficiente y procesos
manuales apoyados en WhatsApp y registros en papel. Esta forma de operar genera
pérdida y dispersión de información, ausencia de trazabilidad en la captación e
instalación, gestión manual y repetitiva de la red, y una excesiva dependencia de
una sola persona, lo que limita la eficiencia operativa, incrementa el riesgo de
errores y restringe la capacidad de crecimiento de la empresa.
\end{upeubox}

El problema puede descomponerse en tres dimensiones interrelacionadas,
correspondientes a los tres flujos operativos críticos de la empresa:

\begin{center}
\captionof{table}{Descomposición del problema por flujo operativo.}
\label{tab:dimensiones}
\begin{tabularx}{\textwidth}{L{2.8cm} X X}
\toprule
\textbf{Dimensión / Flujo} & \textbf{Problema específico} & \textbf{Consecuencia} \\
\midrule
Captación (vendedoras) &
Registro de prospectos en papel o WhatsApp, sin validación de cobertura. &
Pérdida de datos, prospectos en zonas sin fibra, retrabajo administrativo. \\
Instalación (técnicos) &
Evidencias de instalación enviadas por WhatsApp, sin historial técnico. &
Imposibilidad de rastrear la instalación para soporte y mantenimiento futuro. \\
Administración (gestión) &
Alta/corte/reactivación manual en 13 MikroTiks y datos dispersos. &
Cuellos de botella, errores, dependencia de una persona, baja escalabilidad. \\
\bottomrule
\end{tabularx}
\end{center}

\subsection{Análisis de causas raíz}

Aplicando el enfoque de causa--efecto (diagrama de Ishikawa), se agrupan las causas
raíz del problema en seis categorías:

\begin{center}
\captionof{table}{Causas raíz del problema agrupadas por categoría (Ishikawa).}
\label{tab:causas}
\begin{tabularx}{\textwidth}{L{2.6cm} X}
\toprule
\textbf{Categoría} & \textbf{Causas raíz identificadas} \\
\midrule
Tecnología &
Sistema heredado obsoleto (PHP); ausencia de integración con los MikroTiks; sin
aplicación móvil; frontend lento. \\
Procesos &
Procesos de captación, instalación y gestión no digitalizados ni estandarizados;
dependencia de WhatsApp como flujo de trabajo. \\
Personas &
Un solo administrador concentra la operación; personal de campo sin herramientas
adecuadas para su celular. \\
Información / Datos &
Datos dispersos y duplicados; sin fuente única de verdad; información crítica embebida
en comentarios de red o chats. \\
Infraestructura &
Conexión de zonas por SSTP; falta de un canal seguro y unificado (VPN) para que el
sistema central administre los MikroTiks. \\
Gestión / Control &
Sin indicadores ni trazabilidad; cortes y reactivaciones gestionados de forma reactiva
y manual. \\
\bottomrule
\end{tabularx}
\end{center}

\subsection{Impacto estratégico}

El problema descrito no es meramente operativo: compromete directamente el logro de
los objetivos estratégicos de WISP y su sostenibilidad a mediano plazo.

\begin{center}
\captionof{table}{Impacto del problema sobre los objetivos estratégicos.}
\label{tab:impacto}
\begin{tabularx}{\textwidth}{L{3.6cm} X}
\toprule
\textbf{Objetivo estratégico afectado} & \textbf{Impacto del problema} \\
\midrule
OE-01 Expansión &
La carga manual limita la capacidad de crecer a nuevos distritos sin colapsar la
administración. \\
OE-02 Eficiencia operativa &
Los procesos manuales y el retrabajo consumen tiempo y elevan el costo por cliente
gestionado. \\
OE-03 Satisfacción del cliente &
Las demoras en activación y la falta de trazabilidad afectan la experiencia y la
retención. \\
OE-04 Control financiero &
La gestión reactiva de morosidad y la dispersión de datos dificultan el control de
ingresos. \\
OE-05 Escalabilidad tecnológica &
La ausencia de una plataforma integrada impide que la empresa escale de forma
sostenible. \\
\bottomrule
\end{tabularx}
\end{center}

\subsection{Justificación de la intervención}

La intervención propuesta consiste en el desarrollo e implementación de un
\textbf{Sistema de Gestión Integral para ISP de Fibra Óptica Multi-Zona}, compuesto
por un sistema web de administración (Django $+$ Angular $+$ PostgreSQL) desplegado en
la nube, una aplicación móvil con roles de técnico y vendedor, y la integración
directa con los 13 MikroTiks mediante una VPN (WireGuard) y la API de RouterOS, junto
con la automatización de notificaciones vía WhatsApp.

La intervención se justifica desde cuatro perspectivas:

\begin{description}[leftmargin=0pt, style=nextline]
  \item[Justificación operativa:] elimina el uso de WhatsApp y papel como sistema de
  información, digitaliza los tres flujos críticos y reduce drásticamente el retrabajo
  y los errores.
  \item[Justificación estratégica:] habilita la escalabilidad necesaria para expandirse
  a nuevos distritos, alineándose con los objetivos estratégicos de crecimiento y
  eficiencia.
  \item[Justificación económica:] sobre una base estimada de $\approx$455 clientes con
  un ticket promedio aproximado de S/~55 mensuales, la operación gestiona del orden de
  S/~25\,000 mensuales; mejorar su control y reducir la morosidad y el retrabajo tiene un
  retorno directo (detallado en el Business Case).
  \item[Justificación tecnológica:] aprovecha la infraestructura madura existente
  (MikroTik RouterOS 7 con API, fibra GPON) y tecnologías accesibles (cloud, WireGuard)
  para construir una solución moderna y sostenible.
\end{description}

\begin{upeubox}{Conclusión del diagnóstico}
WISP posee una infraestructura de red sólida y una posición competitiva privilegiada
en su mercado, pero su gestión digital se encuentra en un nivel inicial que frena su
crecimiento y su eficiencia. Cerrar esta brecha mediante un sistema de gestión integral
---abordado conjuntamente por las líneas de Gestión de TI, Software e Infraestructura---
es una intervención necesaria, viable y estratégicamente alineada, cuyo detalle de
viabilidad económica se desarrolla en el Business Case (Entregable~2).
\end{upeubox}
